4. Considera el siguiente sistema de ecuaciones.
\[
\begin{array}{l}
x+y+k z=1 \\
x+k y+z=1 \\
k x+y+z=-2
\end{array}
\]

\begin{solution}
Meciona para qué valores de \( k \), si hay alguno, el sistema tiene \\
(a) ninguna solución. \\
(b) una solución. \\
(c) un numero infinito de soluciones. \\

El sistema puede escribirse en forma matricial \( A \mathbf{x} = \mathbf{b} \), donde:
\[
A = \begin{pmatrix}
1 & 1 & k \\
1 & k & 1 \\
k & 1 & 1
\end{pmatrix}, \quad
\mathbf{x} = \begin{pmatrix}
x \\
y \\
z
\end{pmatrix}, \quad
\mathbf{b} = \begin{pmatrix}
1 \\
1 \\
-2
\end{pmatrix}
\]
El determinante de A nos ayudará a determinar la unicidad de la solución mediante Teorema de Rouché-Frobenius.
\[
\det(A) = \begin{vmatrix}
1 & 1 & k \\
1 & k & 1 \\
k & 1 & 1
\end{vmatrix}
\]
Calculamos el determinante usando la regla de Sarrus o la expansión por cofactores:
\[
\det(A) = 1 \cdot (k \cdot 1 - 1 \cdot 1) - 1 \cdot (1 \cdot 1 - k \cdot 1) + k \cdot (1 \cdot 1 - k \cdot k)
\]
\[
\det(A) = 1 \cdot (k - 1) - 1 \cdot (1 - k) + k \cdot (1 - k^2)
\]
\[
\det(A) = (k - 1) - (1 - k) + k (1 - k^2)
\]
\[
\det(A) = k - 1 - 1 + k + k - k^3
\]
\[
\det(A) = 3k - 2 - k^3
\]
\[
\det(A) = -k^3 + 3k - 2
\]
Simplificando:
\[
\det(A) = -k^3 + 3k - 2
\]
Nótese que...\\
Si \( \det(A) \neq 0 \): El sistema tiene una única solución. \\
Si \( \det(A) = 0 \): El sistema puede ser: \\
Si es consistente e indeterminado tiene infinitas soluciones; si es inconsistente implica que no tiene solución. Primero, encontremos los valores de k que hacen que \( \det(A) = 0 \):
\[
-k^3 + 3k - 2 = 0 \quad \Rightarrow \quad k^3 - 3k + 2 = 0
\]
Buscamos las raíces de la ecuación cúbica \( k^3 - 3k + 2 = 0 \). \\
Véase que es factorizable como $ k^3 - 3k + 2 = (k - 1)^2 (k + 2) $. Por lo tanto, las soluciones son:
\[
k = 1 \quad (\text{de multiplicidad 2}), \quad k = -2
\]
Veamos cuidadosamente los siguientes casos: \\
Caso 1: \( \det(A) \neq 0 \) \\
Es decir, cuando \( k \neq 1 \) y \( k \neq -2 \). \\
Conclusión: El sistema tiene una única solución. \\


Caso 2: \( \det(A) = 0 \), es decir, \( k = 1 \) o \( k = -2 \). \\
Analizamos cada valor:


Para $k=1$: \\
Sustituimos $k=1$ en el sistema:
\[
     \begin{cases}
     x + y + z = 1 \\
     x + y + z = 1 \\
     x + y + z = -2
     \end{cases}
     \]
Observamos que las dos primeras ecuaciones son iguales, pero la tercera es inconsistente \( 1 \neq -2 \). \\
Conclusión: No hay solución (sistema inconsistente). \\


Para \( k = -2 \):
Sustituimos $k=−2$ en el sistema:
\[
     \begin{cases}
     x + y - 2z = 1 \\
     x - 2y + z = 1 \\
     -2x + y + z = -2
     \end{cases}
     \]Resolvemos el sistema:
Resolvemos el sistema:
\begin{enumerate}
    \item De la primera ecuación: x=1−y+2z.
    \item Sustituimos x en la segunda y tercera ecuación:
    \item Segunda ecuación:
\[
          (1 - y + 2z) - 2y + z = 1 \quad \Rightarrow \quad 1 - 3y + 3z = 1 \quad \Rightarrow \quad -3y + 3z = 0 \quad \Rightarrow \quad y = z
          \]


    \item Tercera ecuación:
\[
          -2(1 - y + 2z) + y + z = -2 \quad \Rightarrow \quad -2 + 2y - 4z + y + z = -2
          \]
\[
          3y - 3z = 0 \quad \Rightarrow \quad y = z
          \]
    \item De \( y = z \) y sustituyendo en \( x = 1 - y + 2z \):
\[
        x = 1 - z + 2z = 1 + z
        \]
\end{enumerate}
Conclusión: Hay infinitas soluciones parametrizadas por z, de forma elegante este espacio afín es como sigue.
\[
\begin{pmatrix} x \\ y \\ z \end{pmatrix} = \begin{pmatrix} 1 + t \\ t \\ t \end{pmatrix} \; : \; t\in\mathbb{R}  = \begin{pmatrix} 1 \\ 0 \\ 0 \end{pmatrix} + span \Bigg\{ \begin{pmatrix} 1 \\ 1 \\ 1 \end{pmatrix} \Bigg\}
\]
$\therefore$ De forma condensada. \\
(a) Ninguna solución: Para $k=1$. \\
(b) Una única solución: Para cualsea $k\in \mathbb{R}\setminus\{-2,1\}$ \\
(c) Un número infinito de soluciones: Para $k=−2$.
\end{solution}